% \section{Scalable Pre-training of \method}
\section{Model Design and Acceleration}
\label{sec:coding}

\subsection{Spatio-temporal Variational Autoencoder}
\label{sec:vae_arch}

\begin{figure}[th]  
    \centering  
    \includegraphics[width=0.99\textwidth]{./figures/vae/vae.pdf}  
    \caption{Our \method-VAE Framework. \method-VAE can compress the spatio-temporal dimension of a video by $4\times8\times8$ times. The orange rectangles represent $2\times$ spatio-temporal compression, and the green rectangles represent $2\times$ spatial compression.} 
    \label{fig:vae_network} % 
\end{figure}


Variational Autoencoders (VAEs)~\citep{wu2024improved, blattmann2023svd, openaisora2024}  play a crucial role in learning compact latent representations from high-dimensional visual data (especially videos), facilitating scalable and efficient training of generative models, \emph{e.g.}, diffusion models.
%
However, designing effective VAEs for video generation tasks faces several challenges. 
First, videos inherently possess both spatial and temporal dimensions, requiring the VAE to capture complex spatio-temporal dependencies. 
Second, the inherent high-dimensional nature of video (\emph{e.g.}, multiple frames with high-resolution pixels) increases memory consumption and computational costs, making it difficult to scale VAEs to long video sequences.
Third, ensuring temporal causality (\emph{i.e.}, future frames do not influence past frames) is critical for generating realistic and coherent video content, yet this constraint introduces additional architectural complexities. 


To address these challenges, we propose a novel 3D causal VAE architecture (\method-VAE) specifically designed for video generation. 
We combine multiple strategies~\citep{wu2024improved} to improve spatio-temporal compression, reduce memory usage, and ensure temporal causality.
These enhancements make \method-VAE more efficient and scalable and better suited for integration with diffusion-based generative models such as DiT.
Next, we introduce the network design, training details, efficient inference, and experimental comparison of our \method-VAE.


\subsubsection{Model Design}

To achieve a bidirectional mapping between the high-dimensional pixel space and the low-dimensional latent space, we design a 3D causal VAE as illustrated in Fig.~\ref{fig:vae_network}. 
Given an input video $V\in \mathbb{R}^{(1+T) \times H \times W \times 3}$, \method-VAE compresses its spatio-temporal dimensions to $[1+T/4, H/8, W/8]$ while expanding the number of channels $C$ to 16.
Specifically, the first frame is only spatially compressed to better handle the image data, following MagViT-v2~\citep{yu2023language}.
In terms of architecture, we replace all GroupNorm layers~\citep{wu2018group} with RMSNorm layers~\citep{rms} to preserve temporal causality.
This modification enables the use of the feature cache mechanism (refer to Sec.~\ref{sec:cache_infer}), significantly improving the inference efficiency.
Furthermore, we halve the input feature channel in the spatial upsampling layer, resulting in a 33\% reduction in memory consumption during inference.

By carefully tuning the number of base channels further, \method-VAE achieves a compact model size of only 127M parameters. 
This optimization reduces the encoding time and memory usage, thus benefiting the training of subsequent diffusion transformer models.

\subsubsection{Training}

We adopt a three-stage approach to train \method-VAE. 
First, we construct a 2D image VAE with the same structure and train it on image data. Then, we inflate the well-trained 2D image VAE into a 3D causal \method-VAE to provide an initial spatial compression prior, which drastically improves the training speed compared to training a video VAE from scratch. 
At this stage, \method-VAE is trained on the low-resolution (128$\times$128) and small frame number (5-frame) videos to accelerate convergence speed. 
The training losses include $L1$ reconstruction loss, KL loss, and LPIPS perceptual loss, which are weighted by coefficients of 3, 3e-6, and 3, respectively.
Finally, we fine-tune the model on high-quality videos with different resolutions and frame numbers, integrating the GAN loss~\citep{esser2020taming} from a 3D discriminator. 
 

\begin{figure}[t]  
    \centering  
    \includegraphics[width=0.99\textwidth]{./figures/vae/vae_cache.pdf}  
    \caption{Our feature cache mechanism. (a) and (b) show how we use this mechanism in regular causal convolution and temporal downsampling, respectively.} 
    \label{fig:vae_cache} % 
\end{figure}


\subsubsection{Efficient Inference} \label{sec:cache_infer} 


To efficiently support the encoding and decoding of arbitrarily long videos~\citep{CDT,yao2021wenet}, we implement a feature cache mechanism within the causal convolution module of \method-VAE. 

Specifically, the number of video sequence frames follows the $1+T$ input format, so we divide the video into $1+T/4$ chunks, which is consistent with the number of latent features.
When processing input video sequences, the model employs a chunk-wise strategy where each encoding and decoding operation handles only the video chunk corresponding to a single latent representation.
Based on the temporal compression ratio, the number of frames in each processing chunk is limited to at most 4, effectively preventing memory overflow.

To ensure temporal continuity between context chunks, our model strategically maintains frame-level feature caches from preceding chunks.
These cached features are systematically integrated into the causal convolution computations of subsequent chunks.
The specific process is illustrated in Fig.~\ref{fig:vae_cache}, which demonstrates two typical scenarios in our feature cache mechanism. 
In Fig.~\ref{fig:vae_cache} (a), in our default setting, the causal convolution does not change the number of frames, and we need to maintain two cached features (convolution kernel size is 3) from historical frames. 
For the initial chunk, we apply zero padding with two dummy frames to initialize the cache buffer. 
Subsequent chunks systematically reuse the last two frames from the preceding chunk as cached features while discarding obsolete historical data.
Fig.~\ref{fig:vae_cache} (b) presents the scenario involving 2$\times$ temporal downsampling (stride is equal to 2). 
Here, the scenario necessitates different cache management: we implement single-frame cache padding exclusively for non-initial chunks to ensure dimensional consistency. This configuration guarantees that the output sequence length follows the exact downsampling ratio while maintaining causal relationships across the boundaries of the chunk.

This feature cache mechanism not only optimizes memory utilization but also preserves feature coherence across chunk boundaries, thereby supporting stable inference for infinite-length videos.


\subsubsection{Evaluation}  


\begin{figure}[tbp]  
    \centering  
    \includegraphics[width=0.82\textwidth]{./figures/vae/vae_psnr.pdf}  
   \caption{Comparison of video reconstruction performance at $720\times720$ resolution and 25 frames.} 
    \label{fig:vae_psnr}  
\end{figure}

 
\textbf{Quantitative results.}
This study conducts a comprehensive evaluation of the performance of various state-of-the-art (SOTA)  video VAE models~\citep{kong2024hunyuanvideo,cogvideox} by analyzing their Peak Signal-to-Noise Ratio (PSNR) and processing efficiency (frame per unit latency (second)). 
Notably, for fair comparison, most models~\citep{kong2024hunyuanvideo,cogvideox} adopt the same compression rate and latent feature dimension as our approach, \textit{i.e.,} a compression rate of $4\times8\times8$ with a latent dimension of 16.
Open Sora Plan~\citep{yuan_opensora}  has the same compression ratio as ours, but the latent dimension is 4.
SVD~\citep{blattmann2023svd} employs a compression rate of $1\times8\times8$ with a latent dimension of 4,
Step Video~\citep{step_video} employs a compression rate of $8\times16\times16$ with a latent dimension of 64, 
and Mochi~\citep{genmo2024mochi} uses a compression rate of $6\times8\times8$ with a latent dimension of 12.
The evaluation is performed on 200 videos.
These videos consist of 25 frames, and the resolution is set to $720 \times 720$.
In Fig.~\ref{fig:vae_psnr}, the size of the circles is positively correlated with the number of model parameters.
 
Experimental results indicate that \method-VAE exhibits highly competitive performance across both metrics, showcasing a compelling dual advantage of superior video quality and high processing efficiency. 
It is worth noting that under the same hardware environment, our VAE's reconstruction speed is 2.5 times faster than the existing SOTA method (\textit{i.e.}, HunYuan Video). 
This speed advantage will be further demonstrated at higher resolutions due to the small size design of our VAE model and the feature cache mechanism.

This advancement offers an efficient solution for video reconstruction tasks and video generation training. 
These findings not only validate the effectiveness of the proposed model design but also provide valuable insights for the future development of VAE technologies.

\begin{figure}[tbp]  
    \centering  
    \includegraphics[width=0.93\textwidth]{./figures/vae/vae_visual.pdf}  
    \caption{Visualization results of video reconstruction across different scenarios, including texture (first row),  face (second row), text (third row), and high-motion (fourth row). The above videos have rich details but are not used in training.} 
    \label{fig:vae_visual}  
\end{figure}
 

\textbf{Qualitative results.}
 To validate the video reconstruction performance of our method across diverse scenarios, Fig.~\ref{fig:vae_visual}  presents visual results for texture, face, text, and high-motion scenes. 
 Compared to existing VAE models, \method-VAE demonstrates superior performance in these contexts. 
 In the texture scene shown in the first row, our method is capable of capturing details more accurately, such as the texture and direction of hair. 
 In the facial scene depicted in the second row, \method-VAE preserves facial features while reducing blurring and distortion around the lips. 
 In the text scene illustrated in the third row, our method successfully restores textual content with clarity, mitigating issues of character distortion and loss. 
 In the high-motion scene presented in the fourth row, \method-VAE maintains the motion sharpness of video frames. 
 These results indicate that our method offers significant advantages in handling a variety of complex scenarios.
